\documentclass[lang=cn,a4paper,newtx]{cnpc_xsy}

\title{基于AI智能体的经营资料归集与数据分析处理的应用研究}
\author{hsuy1ee \\ 中石油新疆销售有限公司伊犁分公司新源经营部}
\date{\zhdate{2026/04/12}}

\usepackage{siunitx}
\usepackage{xcolor}
\usepackage{listings}

\addbibresource[location=local]{ref.bib}
\renewcommand{\lstlistingname}{示意图}

\begin{document}

\pagenumbering{gobble}

\maketitle
\vspace{1.5em}

\begin{abstract}
    本文围绕 AI 智能体在经营分析资料归集与数据分析处理中的应用展开研究，重点讨论其在办公赋能、流程优化与效率提升方面的实践路径与应用价值。
    结合经营分析资料整理、站点资料归集、考勤核验、工具推广与应用边界等实际场景，文章认为 AI 智能体已经能够在规则明确、重复度高、信息密集的办公任务中形成显著助力，但在高风险现场作业、复杂情境判断和长链条自主执行方面仍存在明显局限。基于此，本文主张以务实方式推进 AI 智能体在基层经营管理中的应用，将其作为提升效率、沉淀规则和促进数智化办公转型的重要工具，同时保持对其能力边界和现实约束的清醒认识。
    \keywords{AI智能体，办公赋能，提升效率，Claude Code，Codex CLI，CodeBuddy}
\end{abstract}

\clearpage
\tableofcontents
\thispagestyle{empty}
\clearpage
\pagenumbering{arabic}
\setcounter{page}{1}

\section{现状与痛点}

\subsection{基层经营管理场景中的重复性资料处理负担}

在片区综合岗位的日常工作中，经营资料的收集、甄别、重命名与归档等基础性事务占据了同事们的部分工作精力。此类事务单次操作复杂度并不高，但频次高、覆盖面广、格式差异大，因而形成了持续存在的隐性工作负担。以月度综合财务检查、节假日值班安排、人员考勤汇总、技能鉴定统计以及库存盘点资料为例，上报文件虽来源不同、命名方式不一，但通常都遵循较为稳定的业务语义结构，即文件名中往往包含站点名称、业务主题和时间属性等关键字段。

例如，综合财务检查资料的文件命名，通常由“站点名称”“财务合规检查问题汇总表”等核心字段，与少量个性化补充表述共同构成；节假日值班资料通常体现为“站点名称+节假日名称+值班表”的结构；考勤等资料也普遍具有“站点名称+业务主题”的共同规律。再特殊一些的库存盘点资料，虽然往往以目录而非单一文件的形式出现，但其内部通常仍包含核对表、盘点表、临期商品台账、异常库存明细等子项，这些文件同样没有脱离“关键字+站点名称”的基本命名逻辑。即使个别文件命名不够规范，相关语义信息仍可从文件标题、正文内容甚至附属元数据中被识别出来。这表明，基层经营资料管理虽然表面上呈现为非标准化问题，但其底层实际上具有较强的规则性、可抽取性与可自动化潜力。

\subsection{经营分析资料处理中的认知负荷问题}

除资料归档外，片区综合岗位还需要定期处理大批量经营分析数据。例如，在月度或季度非油业务分析过程中，工作人员往往需要同时面对商品大类库存明细、收入毛利明细、劳动竞赛数据以及其他经营报表。这些文件来源多元、体量较大，单个文件规模甚至可接近 \SI{50}{\mega\byte}。在连续查阅、筛选和比对大量表格之后，工作人员不仅需要提炼关键指标，还需进一步形成综合性业务简报，其工作过程具有明显的高重复性与高注意力消耗特征。

从认知机理看，这种主观体验并非个体化感受，而具有明确的研究依据。相关研究指出，长时间持续执行单调且高强度的信息处理任务，会引发心理疲劳，并表现为反应时间延长、正确率下降以及对关键信息的注意资源分配减少\cite{guo2016mental,oken2006vigilance}。这与基层经营分析中的实际体验高度一致，即在长时间面对结构相近、指标密集的数据报表时，工作人员容易出现注意稳定性下降、信息辨识效率降低以及短时理解迟滞等现象。由此可见，当前经营分析工作的主要瓶颈不仅来自业务量本身，更来自“重复识别+持续比较+人工归纳”所叠加产生的认知负荷。

\subsection{规则稳定性为智能化介入提供了现实基础}

值得注意的是，此类工作虽然耗时耗力，但并非不可描述、不可建模。无论是资料命名规则，还是月报、季报中反复出现的字段结构、指标口径和分析模板，均具有较强的稳定性。大多数经营分析任务在不同周期之间的差异，主要体现在数据取值、文件命名和局部异常点上，而非业务逻辑的根本变化。因此，这类任务天然适合通过程序化工具或 AI 智能体进行重构。

从流程视角看，若能够将“识别文件类型”“提取关键字段”“按统一规则重命名”“汇总核心指标”“生成分析初稿”等步骤转化为可执行的规则链条，则原本依赖人工经验与反复操作的事务性工作，就可以被迁移为“人制定规则、智能体执行流程、人工复核结果”的新型协同模式。这种模式并不取代业务人员的专业判断，而是将其从低附加值、强重复性的劳动中释放出来，使其更聚焦于异常识别、经营诊断和管理决策支持等更具价值的工作环节。

\subsection{自然语言交互拓展了办公自动化边界}

除文件归集与报表分析外，基层办公场景还广泛存在另一类典型问题，即工具功能可用但使用门槛较高。以 Excel 为例，工作人员在处理数据时常常会出现函数语法遗忘、函数选型不确定或复杂嵌套公式难以构造等情况。传统解决方式通常依赖经验查找、网络搜索或反复试错，效率有限，且容易打断原有工作节奏。

随着大语言模型技术的发展，自然语言到结构化表达的转换能力已在办公场景中展现出明显潜力。Zhao 等提出的 NL2Formula 研究表明，大语言模型及相关方法能够根据自然语言查询生成可执行的电子表格公式，从而为复杂函数编写提供新的技术路径\cite{zhao-etal-2024-nl2formula}。这意味着，基层员工可以逐步从“记忆函数语法”的工作方式，转向“描述业务意图，由模型辅助生成公式”的交互方式。对经营管理岗位而言，这不仅提升了办公软件的使用效率，也进一步说明 AI 智能体在资料归集、数据整理与辅助分析等场景中具有现实可行性和推广价值。

\section{AI智能体在资料归集场景中的应用实践}

\subsection{典型场景一：站点证照与影像资料的自动归集}

在实际工作中，笔者曾处理过一项具有代表性的资料整理任务。该任务面向新源片区共计 19 个站点，原始目录中包含两类文件夹和大量散落照片。其中，两类文件夹分别为不动产证资料和营业执照正本资料；其余则为 19 个站点对应的便利店门头照片、便利店内景照片以及《仅销售预包装食品经营者备案信息采集表》照片。除前述两个目录外，当前目录下其余文件均为 JPG 格式照片，共计 57 张。该批资料在来源、文件命名和存放结构上均存在明显差异，尤其是部分照片命名极不统一，难以通过传统文件检索方式直接完成批量归档。

类似的问题如果完全依靠人工处理，工作人员通常需要反复打开图片或文档，识别分类依据信息，再按照目标结构逐项建立文件夹、复制剪切文件并进行校验。这一过程不仅机械，而且极易因命名相近、文件数量多而发生遗漏或误放。若引入 AI 智能体，则可通过“站点识别+资料类别判断+目标路径映射”的规则体系，实现批量化整理。其结果不是简单地“做了这件事”，而是通过智能体的介入，形成了一个可复用的业务能力模块，能够在未来类似场景中快速复制和推广。

\subsubsection{整理前后的目录结构示意}

为更直观地呈现 AI 智能体介入前后的差异，本文以目录树形式展示原始资料结构与目标归档结构。整理前的特点是“少量目录+大量散落照片并存”；整理后的特点则是“按站点统一归集，同一站点下集中放置相关资料”。

\newpage

\begin{lstlisting}[style=treeview,caption={整理前的示意性目录结构}]
资料原始目录/
├── 不动产证/
│   ├── 卡普河站不动产/
│   ├── 天河站不动产证/
│   └── ...
├── 营业执照正本/
│   ├── 营业执照正本~卡普河站.jpg
│   ├── 天河站营业执照正本.jpg
│   └── ...
├── 便利店门头~卡普河加油站.jpg
├── 便利店内~卡普河加油站.jpg
├── 天河站便利店门头.jpg
├── 天河站便利店内.jpg
├── 仅销售预包装食品经营者备案信息采集表~卡普河站.jpg
├── 天河站仅销售预包装食品经营者备案信息采集表.jpg
└── ...
\end{lstlisting}

\begin{lstlisting}[style=treeview,caption={整理后的示意性目录结构}]
站点资料标准目录/
├── 卡普河站/
│   ├── 不动产证_卡普河站/
│   ├── 营业执照正本_卡普河站.jpg
│   ├── 便利店门头_卡普河站.jpg
│   ├── 便利店内_卡普河站.jpg
│   └── 仅销售预包装食品经营者备案信息采集表_卡普河站.jpg
├── 天河站/
│   ├── 不动产证_天河站/
│   ├── 营业执照正本_天河站.jpg
│   ├── 便利店门头_天河站.jpg
│   ├── 便利店内_天河站.jpg
│   └── 仅销售预包装食品经营者备案信息采集表_天河站.jpg
└── ...
\end{lstlisting}

\subsection{典型场景二：站点考勤表的批量识别、重命名与缺报核验}

除静态资料归集外，AI 智能体在月度经营管理资料的批量整理场景中同样具有较强的适用性。以人员考勤资料为例，站点上报文件往往存在命名不统一、站点简称混用、合同化与劳务派遣材料并存等问题，导致后续汇总、追踪与催报工作需要投入额外精力。针对这一类场景，智能体不仅需要完成重命名，还需要结合站点台账判断哪些站点尚未提交材料，并输出便于后续追踪的缺报清单。

本文选取新源片区 \verb|D:/工作/CNPC/片区/新源片区/综合岗/综合办公室/人员考勤/202603/| 目录作为演示路径，采用命令行方式调用智能体执行任务。演示要求包括：一是将当前目录下的 Excel 文件统一重命名为 \lstinline|"202603_xxx加油站考勤统计表"| 或 \lstinline|"202603_xxx加气站考勤统计表"| 的格式；二是在同一站点同时存在合同化员工与劳务派遣员工考勤表的情况下，确保两个文件不会被错误覆盖，并在必要时通过表格结构或表内内容进一步判断；三是对于 \lstinline|"考勤表（2025版）.xlsx"| 这类文件名无法直接体现站点信息的材料，允许智能体打开 Excel，依据表头、站点名称栏或签章信息识别其实际对应站点；四是读取 \lstinline|"站点+站长.txt"| 中的站点台账，在任务执行完毕后自动生成或覆写一份 \lstinline|"未提交站点.txt"|；五是沉淀一份可复用的 \lstinline|AGENTS.md| 操作规范，并在完成后删除原始 \lstinline|"站点+站长.txt"| 文件。

\begin{codeexample}[htbp]
   \caption{命令行演示与任务指令}
   \begin{lstlisting}[style=cmdview]
D:\工作\CNPC\片区\新源片区\综合岗\综合办公室\人员考勤\202603>codex --yolo

重命名当前目录下的Excel文件，格式为：
"202603_xxx加油站考勤统计表"或"202603_xxx加气站考勤统计表"

要求：
1. xxx在原文件中会有体现
2. 区分同一站点合同化员工和劳务派遣员工，不能导致两个相异的文件被重复命名；
   如若无法区分，可以查看Excel表结构和内容进行判断。
3. "站点+站长.txt"是新源片区所有站点和站长的文件，执行完参考这个文件生成
   或覆写一份"未提交站点.txt"文件以供我查询谁未提交
4. 执行完所有指令之后，生成一份AGENTS.md文件，固定规范复用执行；
   "站点+站长.txt"的内容同样写入md文件；最后删除"站点+站长.txt"
\end{lstlisting}
\end{codeexample}

\subsubsection{整理前后的目录结构示意}

在该场景下，原始目录并非站点维度的结构化资料，而是多个文件名风格并存的散乱目录。通过智能体执行后，文件将按照统一规则完成识别与重命名，并额外生成 \lstinline|"未提交站点.txt"| 用于缺报追踪。为便于对比，本文采用上下顺排方式展示整理前后的目录状态。

\begin{lstlisting}[style=treeview,caption={整理前的示意性目录结构}]
伊犁公司新源县加气站2026年3月合同化员工考勤表.xlsx
伊犁公司阿勒玛勒加油站2026年3月合同化员工考勤表（2025年新版).xlsx
副本铁新公路站-新源片区2026年3月考勤统计表（202408新版）.xlsx
新源县加气站2026年3月劳务派遣考勤表.xlsx
新源片区阿勒玛勒2026年3月考勤统计表（202408新版）.xlsx
木斯服务区北加油站站2026年3月考勤统计表（202408新版）.xlsx
木斯服务区南加油站站2026年3月考勤统计表（202408新版）.xlsx
种羊场加气站员工考勤表-2026年3月.xlsx
种羊场加油站2026年3月考勤统计表（202408新版）.xlsx
考勤表（2025版）.xlsx
那拉提高速北加油站2026年3月考勤统计表.xlsx
那拉提高速南加油站2026年3月考勤统计表-.xlsx
站点+站长.txt
\end{lstlisting}

\newpage

\begin{lstlisting}[style=treeview,caption={整理后的示意性目录结构}]
202603_卡普河站考勤统计表.xlsx
202603_阿勒玛勒站合同化员工考勤统计表.xlsx
202603_阿勒玛勒站劳务派遣员工考勤统计表.xlsx
202603_铁新加油站考勤统计表.xlsx
202603_新源县加气站合同化员工考勤统计表.xlsx
202603_新源县加气站劳务派遣员工考勤统计表.xlsx
202603_木斯服务区北站考勤统计表.xlsx
202603_木斯服务区南站考勤统计表.xlsx
202603_种羊场加气站考勤统计表.xlsx
202603_种羊场站考勤统计表.xlsx
202603_那拉提服务区北站考勤统计表.xlsx
202603_那拉提服务区南站考勤统计表.xlsx
未提交站点.txt
AGENTS.md
\end{lstlisting}

\subsubsection{执行结果与缺报核验}

从执行结果看，智能体在该场景中完成了三项关键任务。其一，对文件名中 \lstinline|"新源片区"|、\lstinline|"副本"|、\lstinline|"202408新版"|、\lstinline|"2025版"| 等非核心表述进行了剥离，保留了月份、站点和资料主题等关键字段，使文件命名回归统一规则。其中，原始文件 \lstinline|"考勤表（2025版）.xlsx"| 并未在文件名中直接给出站点名称，智能体需要实际读取 Excel 内容，并根据表内出现的“卡普河加油站”等关键信息判断其对应站点，最终将其规范化为 \lstinline|"202603_卡普河站考勤统计表.xlsx"|。其二，对同一站点的合同化员工考勤表与劳务派遣员工考勤表进行了区分处理，避免出现因同名覆盖导致的数据丢失。其三，结合 \lstinline|"站点+站长.txt"| 中的完整站点台账，生成了 \lstinline|"未提交站点.txt"|，将缺报情况从人工逐一核对转化为自动比对输出。

按照新源片区站点台账，共有 19 个站点纳入比对范围：卡普河站、阿热勒托别站、城北站、则克台镇西站、塔勒德镇西站、则克台镇东站、阿勒玛勒站、哈拉布拉站、种羊场站、天河站、木斯站、新河站、木斯服务区南站、木斯服务区北站、那拉提服务区南站、那拉提服务区北站、铁新加油站、新源县加气站、种羊场加气站。结合演示目录中的已提交材料，可以判定本次未提交站点包括：阿热勒托别站、城北站、则克台镇西站、塔勒德镇西站、则克台镇东站、哈拉布拉站、天河站、新河站、木斯站。由此可见，AI 智能体并非只是在“把文件名改整齐”，而是能够进一步完成站点映射、材料去重、异常识别和缺报输出等成套动作，体现出较强的执行稳定性与业务可靠性。

\subsection{AI智能体驱动的业务价值}

从业务效果看，该类应用至少带来三方面收益。其一，显著降低资料归集所需的重复劳动时间，使工作人员从机械性整理中释放出来。其二，通过统一命名与目录结构，提升资料管理的一致性和可检索性，减少因个人习惯差异造成的后续沟通成本。其三，为经营资料进一步进入结构化管理、知识沉淀和智能分析创造前提条件。换言之，AI 智能体在此类场景中的意义，不仅在于“替人做事”，更在于把原本松散的文件处理过程，转化为可复用、可复制、可扩展的标准化业务能力。

因此，在基层经营管理实践中，AI 智能体并不是附加性的技术装饰，而是推动经营资料管理从经验驱动走向规则驱动、从人工分散走向流程集成的重要抓手。围绕资料归集、数据清洗、报表分析和文本生成等环节持续推进智能化建设，有望形成可推广的数智化办公范式。

\section{AI智能体工具的推广使用路径}

前两节讨论了 AI 智能体在资料归集、文件重命名、缺报核验等场景中的具体作用，但若要将这类能力从个体尝试转化为可复制的办公方法，还需要回答一个更现实的问题，即基层岗位应当如何选择、安装、配置并稳定使用相关工具。与传统办公软件不同，AI 智能体往往依赖命令行环境、模型接口配置以及一定程度的项目规则约束，因此其推广不应停留在“知道有这个工具”，而应形成一套可操作、可迁移、可维护的使用路径。

需要说明的是，本文这里讨论的工具，并不意味着它们构成当前 AI 智能体工具的主流全貌，而只是笔者在实际接触过程中认为较有代表性工具。

\subsection{Claude Code 的安装与 GLM 模型平替配置}

在笔者的实际接触中，Claude Code 是较早进入使用视野的一类工具。其特点在于交互方式较清晰，规则文件思路也比较适合形成固定工作流。如果将其用于资料归集、批量处理、规则复用等任务，较容易从一次性的尝试逐步沉淀为可重复执行的方法。

从安装角度看，Claude Code 依赖 Node.js 与 npm 环境，因此在使用前需要先补齐 Node.js，再执行安装命令。本文这里给出的安装和 GLM 模型平替配置方式，除参考相关文档外，也部分参考了博主：\lstinline|"杰森的效率工坊"| 的视频内容。\cite{youtube_obsidian_claude_skill_2024,zai_claude_code_2026}。

\begin{codeexample}[htbp]
   \caption{Claude Code 安装与环境变量配置示例}
   \begin{lstlisting}[style=cmdview]
# 安装 Claude Code（需先具备 Node.js 与 npm 环境）
npm install -g @anthropic-ai/claude-code

# Windows 环境变量配置示例
setx ANTHROPIC_BASE_URL "https://open.bigmodel.cn/api/anthropic"
setx ANTHROPIC_AUTH_TOKEN "你的_GLM_API_KEY"
setx ANTHROPIC_MODEL "glm-4.7"

# 进入工作目录后启动
cd D:\工作\CNPC\片区\新源片区\综合岗
claude
\end{lstlisting}
\end{codeexample}

\subsection{CodeBuddy 的安装与使用}

除 Claude Code 外，CodeBuddy 也是值得纳入推广讨论的一类工具。根据腾讯云官方资料，CodeBuddy 由腾讯推出，提供 AI 对话、代码补全、工程理解、智能评审、单元测试和 MCP 等能力，并支持多种 IDE 与命令行工作方式\cite{codebuddy_product_2026}。对于国内用户而言，这类工具的现实优势在于本土化服务链路相对完整，安全系数更高，账号获取以及网络可达性往往更容易落地。

CodeBuddy CLI 同样支持通过 npm 安装。根据其快速入门文档，Windows、macOS 和 Linux 环境均可使用，且推荐在安装完成后进入目标项目目录启动工具\cite{codebuddy_quickstart_2026}。这说明它不仅适用于代码开发场景，也可以迁移到规则明确、目录清晰的经营分析资料处理工作中。若使用者不偏好 CLI（Command Line Interface，命令行界面）方式，也可以直接下载 CodeBuddy IDE\footnote{\url{https://www.codebuddy.ai/ide}}。相较而言，CLI 更适合执行连续指令、批量任务和复杂流程，整体能力通常更灵活；而 IDE 更强调图形化界面与交互友好性，更便于初学者理解和上手。

\begin{codeexample}[htbp]
   \caption{CodeBuddy 的基础安装与启动示例}
   \begin{lstlisting}[style=cmdview]
# 安装 CodeBuddy Code
npm install -g @tencent-ai/codebuddy-code

# 进入项目目录并启动
cd D:\工作\CNPC\片区\新源片区\综合岗
codebuddy
\end{lstlisting}
\end{codeexample}

相较于需要额外兼容接口配置的使用方式，CodeBuddy 的推广障碍通常更低一些，更适合作为普通员工的入门型选择。在单位内部开展小范围试点时，这类工具也更便于统一安装说明、统一演示流程和统一答疑口径。当然，这并不意味着应只保留一种工具，而是说明在过渡阶段，可以根据不同岗位条件形成多工具并行、但规则模板相对统一的使用格局。

总体来看，AI 智能体工具的推广，本质上并不只是安装哪个软件的问题，而是如何将零散经验沉淀为共享方法的问题。对基层经营管理岗位而言，真正值得复制的不是某一条安装命令，而是围绕这些工具逐步形成的目录约定、提示词模板、输出格式、校验流程和规则文件。只有当这些方法沉淀下来，AI 智能体才可能从个人效率工具，进一步转化为可复制的数智化办公能力。

\section{AI智能体的不足}

尽管前文已经展示了 AI 智能体在资料归集、数据整理和规则化处理中的显著优势，但这并不意味着其能力可以无条件外推到所有业务环节。特别是在安全要求高、环境变化快、需要持续情境判断和即时协同处置的现场作业中，AI 智能体目前仍存在明显边界。

\subsection{现场巡检与生产作业并非纯信息处理任务}

笔者在 2025 年底于油库工作期间，曾听到关于“以机器人替代人工巡检”的讨论。这一设想在技术演示层面确实具有吸引力，但若据此推断当前 AI 或机器人系统已经能够直接替代油库巡检岗位，显然仍为时过早。以油库巡检为例，完整路线一次最快也需要约 \SI{30}{\minute}，行走距离超过\SI{2}{\kilo\meter}。在这一过程中，工作人员并不是简单地沿既定路线行走，而是需要持续结合气味、温度、声音、震动、仪表状态、地面积液、设备细微变化以及现场作业上下文进行综合判断，并在发现异常时立即与中控或相关岗位联动处置。

以宇树科技的相关产品资料为例，其官方说明已经展示出较强的移动与巡检潜力，例如可通过激光雷达建图、按自定义路网进行自主巡逻与导航，并实时显示状态、温度和报警等信息\cite{unitree_b2_app_2025}。这说明当前机器人已经能够承担一部分“按路线移动+采集部分状态信息”的工作。然而，这与真正意义上的现场安全判断仍存在明显差距。前者更接近预设路径下的移动执行与状态回传，后者则需要面向开放环境进行连续的多模态理解、异常识别、风险分级以及处置决策。

同样地，春晚或技术大会上展示的人形机器人、娱乐机器人虽然令人印象深刻，但其“惊艳”主要体现于动作控制、交互表现、导航能力或策略训练进展，而不能直接等同于已具备在高风险工业环境中独立分析现实情境并承担安全责任的能力。以 NVIDIA 在 GTC 2026 期间发布的机器人生态进展为例，其官方表述强调的是 physical AI、世界模型、仿真训练、导航复杂环境与机器人策略学习等能力建设。例如 Disney 的 Olaf 和 BDX Droids 被用于训练策略并实现复杂环境导航\cite{nvidia_physical_ai_2026}。这说明机器人正朝着更强的感知与行动能力演进，但从“可演示、可训练、可导航”走向“可在危化品场景中稳定替代人工判断与处置”，中间仍有很长距离。

我还想补充：至少在油库场景中，AI 智能体至少暂时无法替代卸油栈桥上依靠经验判断现场状态的卸油人员，也难以替代泵房内持续拿着对讲机、与中控实时沟通并即时控制阀门的老师傅们。这些岗位的核心能力并不是单一动作执行，而是把经验、责任、风险意识和现场变化整合为即时判断。换言之，越是依赖复杂情境感知和责任闭环的现场岗位，越不能简单用当前的智能体能力去类比替代。

\subsection{岗位替代并非线性发生，现实影响仍在缓慢显现}

围绕 AI 是否会迅速替代大量岗位，近年来外部讨论往往较为激进。但从现有研究看，现实中的劳动替代效应并未以一种简单、同步、立刻显现的方式发生。Anthropic 于 2026 年发布的研究指出，AI 对劳动市场的影响需要区分“理论上能够完成哪些任务”和“现实中实际覆盖了多少任务”。该研究发现，AI 的实际覆盖范围仍明显低于其理论能力边界；同时，自 2022 年末以来，并未观察到高暴露岗位出现系统性的失业率上升，只是对年轻劳动者进入高暴露岗位的招聘放缓存在一定迹象\cite{anthropic_labor_market_2026}。

\begin{figure}[htbp]
    \centering
    \includegraphics[width=0.82\linewidth]{anthropic_labor-market-impacts.jpg}
    \caption{Anthropic 关于 AI 劳动市场影响研究中的相关图示\cite{anthropic_labor_market_2026}}
    \label{fig:anthropic-labor-market}
\end{figure}

这类研究结论对基层岗位判断具有重要启发。一方面，AI 的确正在改变部分任务的完成方式，尤其是文书、分析、客服、数据录入和编码等高信息密度场景；但另一方面，真正的岗位替代并不只是“模型能不能做”，还取决于制度责任、现场风险、组织流程、复核要求和软硬件协同条件。也就是说，AI 更可能先改变岗位内部的任务结构，而不是立刻整体消灭某个岗位。

\subsection{AI Agent 自身仍处于不够成熟的阶段}

除了物理现场作业难以完全替代之外，AI agent 自身的发展阶段也决定了它尚不适合被无条件神化。以软件工程领域的严格评测为例，SWE-Bench Pro 专门针对更接近真实企业环境、长时程、多文件修改的软件工程任务进行测试。结果表明，即便是当前领先模型，在该基准上的通过率仍低于 \SI{25}{\percent}，说明现阶段的 agent 在长链条任务、复杂上下文管理和企业级真实问题求解方面，距离稳定、鲁棒、可无人监督部署的状态还有较大差距\cite{deng2025swebenchpro}。

从产业观察角度看，Gartner 在 2025 年的预测也给出了类似提醒：到 2027 年底，超过 \SI{40}{\percent} 的 agentic AI 项目将被取消，其原因主要包括成本上升、业务价值不清晰以及风险控制不足；同时，Gartner 还指出，当前许多 agentic AI 项目仍停留在早期实验或概念验证阶段，且相当部分受到市场热度驱动\cite{gartner2025_agentic_cancelled}。这一判断说明，当前 agent 热潮中既包含真实进步，也夹杂着大量概念放大和场景误配。

因此，若回到基层经营管理实践，AI 智能体更适合被理解为“在特定规则和边界内显著提升效率的工具”，而不是能够脱离监管、脱离复核、脱离现场责任独立完成全部工作的替代者。对资料整理、表格分析、初稿生成、规则匹配这类任务，AI agent 已经具有较高价值；但对长时程、多环节、强责任、强协同、强现场感知的任务，其不足同样必须被正视。只有同时看到能力与边界，AI 智能体才能真正以务实方式服务于经营管理，而不至于陷入“技术神话”。

\section{未来与展望}

写到最后，笔者更愿意把 AI 看作一种正在展开的新风口，而不是一个已经尘埃落定、边界分明的成熟工具体系。也正因为它还远没有形成完全清晰的规矩、标准和使用边界，所以才给了普通人更多尝试的空间。今天的很多用法，也许在半年前还没有被认真讨论；而现在看起来很自然的功能，可能也是很多人不断试出来的结果。从这个意义上说，AI 的发展阶段并不是一种缺陷，反而像是一块仍在生长中的土壤，谁愿意先动手、先学习、先试错，谁就更有机会在其中形成自己的方法和优势。

如果把时间线再拉长一些来看，这种变化其实并不陌生。十多年前，移动互联网兴起的时候，很多人也未必能准确判断它会如何改变生活和工作，但后来大家都身处其中，很难真正置身事外。今天的 AI，某种程度上也正在经历类似的过程。无论从事什么行业，人们都可能逐步发现，自己越来越难脱离 AI 工具而独立完成全部工作。它也许不会在一夜之间改变一切，但会一点一点渗透进资料整理、沟通协同、分析判断、知识获取和内容生产等各个环节，最后变成一种默认存在。

对普通人而言，这里面最重要的意义在于一种相对平等的创造机会。过去很多能力门槛较高的工作，需要较长时间的专业训练才能进入；而现在，借助 AI，越来越多的人可以把自己的想法更快地转化为文字、表格、脚本、流程和方案。一个人的专业背景当然仍然重要，但持续学习、主动尝试和善于提问，正在变得越来越重要。谁能更快地理解工具、组织任务、形成方法，谁就更可能在新环境中建立起新的竞争力。

写到这里，笔者也想起自己刚入职时在昌吉参加培训的经历。当时伊犁公司吴东先书记也来到现场，大家都在向各自地州的前辈请教困惑。笔者当时提过一个问题，大意是：在新的工作环境里，应该怎样持续提升自己的竞争力。吴书记给出的回答，核心意思其实很朴素，就是要始终保持不断学习的动力。现在回头看，这句话依然很适用于今天面对 AI 的我们。工具会不断变化，平台会不断更替，模型也会不断更新，但愿不愿意持续学习、愿不愿意主动适应，始终是决定一个人能走多远的关键。

因此，本文对未来的判断并不复杂：AI 智能体不会取代所有人，但会越来越深地影响每个人的工作方式；不会让学习失去意义，反而会让持续学习变得更加重要；不会把普通人排除在外，反而可能为愿意尝试的人提供新的机会。对于基层经营管理岗位来说，最现实的选择也许不是犹豫“要不要接触 AI”，而是尽早去理解它、使用它、驯化它，让它真正成为提升效率、拓展能力和适应时代变化的一部分。与其等待边界被完全定义，不如在实践中逐步建立自己的用法。拥抱 AI，本质上并不是追逐口号，而是尽量跟上时代向前走的速度。


\printbibliography[heading=bibintoc,title=\ebibname]

\end{document}
